% 2026-11-14 / A-v61.
\section{満supportの同時量子化と全階数のスカラーgerbe}
\label{sec:pg-full-support-quantization}

本節では、正規形に現れるブロックの順序を忘れた係数空間を扱う。
まず全ての非零階数層について、中心スカラーgerbeの位数を求める。
次に$N=3k$、すなわち中心変数が残らない満support層では、
係数空間がアフィンであることを用いて、指定したHessian二次jetを持つ
共変な代数的中心元を同時に選ぶ。
その三次Veroneseは実際のPoisson対の被約モジュライへ降下する。
ブロックの内在性と相対階数層の記述は前節までの結果を使い、
アフィン商・還元群の平均・重み零の降下を既知の入力として区別する。

\subsection{内在的なブロック被覆と全階数層のgerbe}

$N\ge5$、$1\le k\le\lfloor N/3\rfloor$、$m=N-3k$とし、
$V=\CC^N$を固定する。基準となる分解と三重ベクトルを
\begin{equation}
 V=W_1\oplus\cdots\oplus W_k\oplus U,\qquad
 \dim W_s=3,\quad\dim U=m,\qquad
 P_k=\sum_{s=1}^kP_s,\quad
 0\ne P_s\in\bigwedge^3W_s
 \label{eq:pg-intrinsic-multiblock-reference}
\end{equation}
とおく。$U$は以下の表示のために選ぶ補空間であり、
$P_k$から内在的に選ばれているとは主張しない。

\begin{lemma}[安定化群と順序付きブロック被覆]
\label{lem:pg-intrinsic-block-cover}
$W=\operatorname{supp}P_k=\bigoplus_sW_s$とおく。
順序を忘れた集合$\{W_1,\ldots,W_k\}$と各成分$P_s$は、
補題\ref{ne:blocks}により$P_k$から復元される。
$J_{N,k}=\operatorname{Stab}_{\mathrm{GL}(V)}(P_k)$は、
選んだ補空間$U$を用いると
\begin{equation}
 \begin{aligned}
 J_{N,k}
  &=\operatorname{Hom}(U,W)\rtimes
     \bigl((\mathrm{SL}_3^k\rtimes\mathfrak S_k)
                   \times\mathrm{GL}(U)\bigr),\\
 J_{N,k}^{\circ}
  &=\operatorname{Hom}(U,W)\rtimes
     \bigl(\mathrm{SL}_3^k\times\mathrm{GL}(U)\bigr),
 \qquad J_{N,k}/J_{N,k}^{\circ}\simeq\mathfrak S_k
 \end{aligned}
 \label{eq:pg-intrinsic-multiblock-stabilizer}
\end{equation}
となる。ここで各$\mathrm{SL}_3$は$P_s$を保つ群であり、
$\mathfrak S_k$はこの体積形式付きブロックを置換する。
$\mathcal O_{N,k}=\mathrm{GL}(V)/J_{N,k}$上には、
ブロックの集合を表す内在的な次数$k$の有限\'etale被覆がある。
その全順序を選ぶ被覆は
\begin{equation}
 \mathcal O_{N,k}^{\mathrm{ord}}
   =\mathrm{GL}(V)/J_{N,k}^{\circ}
   \longrightarrow\mathcal O_{N,k}
 \label{eq:pg-intrinsic-ordered-block-cover}
\end{equation}
という$\mathfrak S_k$-torsorであり、次数は$k!$である。
これらの構成は底変更と可換である。
\end{lemma}
\begin{proof}
$W$は縮約写像$\bigwedge^2V^*\to V$の像であり、$P_k$から決まる。
$W^*=\bigoplus_sW_s^*$の座標を用いると、
\begin{equation}
 D_{P_k}
  =\{\alpha\in W^*:(\iota_\alpha P_k)\wedge
                    (\iota_\alpha P_k)=0\}
  =\bigcup_{s=1}^kW_s^*
 \label{eq:pg-intrinsic-block-contraction-locus}
\end{equation}
である。ここで右辺は被約なスキームとしての等式である。
実際、左辺の定義イデアルは、異なるブロックの座標の積で生成される。
これは右辺の座標線形空間のイデアルの共通部分である。
これは補題\ref{ne:blocks}の証明そのものであり、
$k=1$では零イデアル、$k=2$でも同じ議論が成立する。
従って$P_k$を保つ変換は$W$を保ち、その制限は$W_s$を置換する。

この変換は$W\oplus U$に関して
\begin{equation}
 L=\begin{pmatrix}A&C\\0&D\end{pmatrix},\qquad
 C\in\operatorname{Hom}(U,W),\quad D\in\mathrm{GL}(U)
 \label{eq:pg-intrinsic-block-stabilizer-matrix}
\end{equation}
と書ける。$A(W_s)=W_{\sigma(s)}$なら、
$\bigwedge^3A(P_s)=P_{\sigma(s)}$が必要十分である。
従って対応する三次元の体積形式に関する行列式は一である。
$C,D$への追加条件はなく、これで
\eqref{eq:pg-intrinsic-multiblock-stabilizer}を得る。
標数零では安定化群は滑らかであり、この表示は閉点だけでなく
代数群の表示である。特に$k=2$にも追加の連続対称性はない。

軌道上では縮約写像の階数が一定なので、その像は階数$3k$の
普遍部分束$\mathcal W\subset V\otimes\OO_{\mathcal O_{N,k}}$となる。
\eqref{eq:pg-intrinsic-block-contraction-locus}を相対化し、
\begin{equation}
 p:\mathbf P(D_P)\longrightarrow\mathcal O_{N,k},\qquad
 \mathcal O_{N,k}^{\mathrm{blk}}
   =\operatorname{Spec}_{\mathcal O_{N,k}}p_*\OO_{\mathbf P(D_P)}
 \label{eq:pg-intrinsic-unordered-block-cover}
\end{equation}
とおく。$\mathrm{GL}(V)\to\mathcal O_{N,k}$で引き戻すと
$\mathbf P(D_P)$は$k$個の$\mathbf P^2$の非交和となる。
従って$\mathcal O_{N,k}^{\mathrm{blk}}$は次数$k$の有限\'etale被覆であり、
その順序付けの$\mathfrak S_k$-torsorが
\eqref{eq:pg-intrinsic-ordered-block-cover}である。
各双対成分から、他の成分の和の零化空間として$W_s$を復元できる。
また$P_s$は$P$の$\bigwedge^3W_s$成分である。
以上は忠実平坦被覆上での構成と降下なので、任意底変更と可換である。
\end{proof}

三次式が$P$と適合するとは、異なる支持ブロックの方向に関する
混合二階偏微分が零であることとする。
これは補題\ref{lem:pg-linearize}の条件である。
基準点$P_k$における適合三次式の空間を$\mathcal F_{N,k}$と書くと、
定理\ref{thm:pg-fermat-relative-count}第三項の証明と同じく
\begin{equation}
 \begin{aligned}
 F&=c(u)+\sum_{s=1}^k\sum_{j=1}^3 f_{s,j}(x_s,u),\\
 \deg_{x_s}f_{s,j}&=j,\qquad\deg_u f_{s,j}=3-j,\\
 \ell_{N,k}:=\dim\mathcal F_{N,k}
   &=\binom{m+2}{3}
     +k\left(10+6m+3\binom{m+1}{2}\right)
 \end{aligned}
 \label{eq:pg-intrinsic-multiblock-cubic-fiber}
\end{equation}
である。その滑らかな三次式の部分を$\mathcal F_{N,k}^{\mathrm{sm}}$とし、
\begin{equation}
 \mathcal S_{N,k}
  =\mathrm{GL}(V)\times^{J_{N,k}}\mathcal F_{N,k}^{\mathrm{sm}}
 \label{eq:pg-intrinsic-multiblock-incidence}
\end{equation}
とおく。これは$\mathcal O_{N,k}$上の階数$\ell_{N,k}$のベクトル束の
判別式非消滅開集合であり、滑らかかつ既約である。
以下ではこの被約なスキーム構造を用いる。
Fermat三次式により非空であることも分かる。
この記述には、純中心三次式$c(u)$をブロックに配分する選択は含まれない。
補題\ref{lem:pg-intrinsic-block-cover}の被覆を引き戻すと、
$\mathcal S_{N,k}$上にも内在的なブロック被覆と順序付け被覆を得る。

\begin{proposition}[全非零階数層の商とgerbeの正確な位数]
\label{prop:pg-intrinsic-all-rank-gerbe}
$\mathfrak P_{n,k}^{\mathrm{red}}$を、$n=N-2$次元の滑らかな三次超曲面と
一般階数$2k$のPoisson構造の相対スタックの被約層とする。
群と作用を
\begin{equation}
 \begin{aligned}
 H_N&=\mathrm{GL}(V)\times\mathbf G_{m,\lambda},\\
 (L,\lambda)\cdot(F,P)
   &=\bigl(\lambda^{-1}F\circ L^{-1},
            \lambda(\bigwedge^3L)P\bigr),\\
 K_{\mathrm{sc}}&=\{(cI_V,c^{-3}):c\in\mathbf G_m\},
 \qquad\overline H_N=H_N/K_{\mathrm{sc}}
 \end{aligned}
 \label{eq:pg-intrinsic-all-rank-action}
\end{equation}
とおく。このとき
\begin{equation}
 [\mathcal S_{N,k}/\overline H_N]
    \simeq\mathfrak P_{n,k}^{\mathrm{red}},\qquad
 \mathfrak G_{N,k}:=[\mathcal S_{N,k}/H_N]
    \longrightarrow\mathfrak P_{n,k}^{\mathrm{red}}
 \label{eq:pg-intrinsic-all-rank-quotient}
\end{equation}
の第二の射は$\mathbf G_m$-gerbeである。
その類を$\alpha_{N,k}$とすると
\begin{equation}
 \begin{aligned}
 \operatorname{im}\bigl(\operatorname{wt}:\operatorname{Pic}(\mathfrak G_{N,k})
                                  \longrightarrow\mathbb Z\bigr)
     &=\gcd(N,3)\mathbb Z,\\
 \operatorname{ord}(\alpha_{N,k})
     &=\gcd(N,3)=\gcd(m,3)
 \end{aligned}
 \label{eq:pg-intrinsic-all-rank-period}
\end{equation}
である。特に$3\mid N$ならこのgerbeは非中立であり、
スカラー重み一の可逆層は存在しない。
$\gcd(N,3)=1$なら中立である。
ここで$\operatorname{wt}$は$K_{\mathrm{sc}}\simeq\mathbf G_m$の作用の重みであり、
$\gcd(0,3)=3$と約束する。
同じ位数は、全零切断を含む
$\mathfrak P_{n,k}^{\mathrm{red}}\times\mathbb A^1_\hbar$の任意の開近傍への
このgerbeの引戻しでも保たれる。
位数の主張は階数層全体、またはその全零切断を含む開近傍についてのものであり、
任意の稠密開集合や一般点における位数を主張しない。
特に満supportの場合$N=3k$には位数は常に三である。
\end{proposition}
\begin{proof}
正規形定理\ref{thm:pg-normal}および定理\ref{thm:pg-fermat-relative-count}第三項の証明により、
三次式の代表を選んだ階数$2k$の被約層は
\eqref{eq:pg-intrinsic-multiblock-incidence}である。
ここで被約構造を指定しており、Jacobi方程式で定まるPoissonスキーム全体の
被約性を仮定しているのではない。
方程式の代表を変える作用$(F,P)\mapsto(\lambda^{-1}F,\lambda P)$は自由であり、
\begin{equation}
 1\longrightarrow\mathbf G_{m,\lambda}
 \longrightarrow\overline H_N
 \longrightarrow\mathrm{PGL}(V)\longrightarrow1
 \label{eq:pg-intrinsic-all-rank-effective-group}
\end{equation}
は完全である。従って代表のスカラーを先に除き、次に$\mathrm{PGL}(V)$で割ると、
\eqref{eq:pg-intrinsic-all-rank-quotient}の第一の同型を得る。
$K_{\mathrm{sc}}$は底に自明に作用する中心部分群なので、第二の射はgerbeである。
ブロック被覆は$P$の非零スカラー倍で不変であり、これらの商にも降下する。

指標$(\det L)^r\lambda^s$の$K_{\mathrm{sc}}$上の重みは$Nr-3s$である。
特に$\lambda^{-1}$は重み三を持つ。
$\gcd(N,3)=1$なら$Nr-3s=1$を選べるので、
$\theta=(\det L)^r\lambda^s$は重み一の指標となる。
この場合は
\begin{equation}
 H_N\simeq K_{\mathrm{sc}}\times\ker\theta,
 \qquad\ker\theta\simeq\overline H_N
 \label{eq:pg-intrinsic-all-rank-neutralization}
\end{equation}
と分裂し、gerbeは中立である。

$3\mid N$の場合の逆向きの制約を示す。
$\zeta$を一の原始三乗根とし、$V$の基底を$N/3$個の三つ組に分ける。
最初の$k$個をactiveブロックとし、
\begin{equation}
 F_{\mathrm F}=\sum_{i=1}^N x_i^3,\qquad
 P_k=\sum_{s=1}^k
 e_{3s-2}\wedge e_{3s-1}\wedge e_{3s}
 \label{eq:pg-intrinsic-period-fermat-point}
\end{equation}
を考える。inactive部分の次元$m=N-3k$も三の倍数なので、
全ての三つ組に同じ変換を施せる。
各三つ組の基底を$e_0,e_1,e_2$と書き、添字を三で読むとき
\begin{equation}
 S(e_j)=e_{j+1},\qquad T(e_j)=\zeta^j e_j,
 \qquad STS^{-1}T^{-1}=\zeta^{-1}I_V
 \label{eq:pg-intrinsic-period-heisenberg-commutator}
\end{equation}
とおく。各三つ組で$\det S=\det T=1$であり、
$S,T$はいずれもFermat三次式と各active体積形式を保つ。
従って$(S,1),(T,1)$は$(F_{\mathrm F},P_k)$の$H_N$-安定化群に属し、
その交換子は
\begin{equation}
 (\zeta^{-1}I_V,1)
   =(cI_V,c^{-3})\big|_{c=\zeta^{-1}}\in K_{\mathrm{sc}}
 \label{eq:pg-intrinsic-period-scalar-commutator}
\end{equation}
である。

$\mathfrak G_{N,k}$上の任意の可逆層は、
$\mathcal S_{N,k}$上の$H_N$-同変可逆層である。
その$(F_{\mathrm F},P_k)$におけるファイバーは安定化群の一次元表現であり、
交換子には自明に作用する。
従ってスカラー重みを$a$とすると$\zeta^{-a}=1$、すなわち$3\mid a$である。
$\mathcal S_{N,k}$は既約なので、この重みは底全体で同じ整数である。
重み三の指標が既にあることから、重みの像はちょうど$3\mathbb Z$となる。
この議論は可逆層の下部の自明性やPicard群の計算を必要としない。

式\eqref{eq:pg-intrinsic-gerbe-weight-period-criterion}とそのcocycle証明により、
代数スタック上の$\mathbf G_m$-gerbeの類$\alpha$に対しても、
重み$a$の可逆層の存在は$a\alpha=0$と同値である。
代数空間の場合の既知の記述は\cite[Lemma~3.1.1.8]{PLieblichPeriodIndex}を参照する。
以上により\eqref{eq:pg-intrinsic-all-rank-period}を得る。
全零切断を含む開近傍では、重み三の可逆層を引き戻せるので位数は三を割り、
零切断への制限により元の位数を下回らない。
$\gcd(N,3)=1$の場合の中立化も引き戻せるので、最後の主張が従う。
\end{proof}

\begin{remark}[満supportの場合のアフィン性]
\label{rem:pg-intrinsic-full-support-affinity}
$m=0$の場合、$J_{N,k}=\mathrm{SL}_3^k\rtimes\mathfrak S_k$は還元的である。
還元部分群による斉次空間のアフィン性\cite[Proposition~2.2, p.~2]{PArzhantsevMatsushima}により、
$\mathcal O_{N,k}$はアフィンである。
従ってその上のベクトル束の全空間もアフィンであり、
判別式非消滅の主開集合$\mathcal S_{N,k}$もアフィンである。

$m>0$では$J_{N,k}$のユニポテント根基
$\operatorname{Hom}(U,W)$は非自明であり、軌道$\mathcal O_{N,k}$自体は
Matsushimaの判定法\cite[Theorem~2.1, p.~2]{PArzhantsevMatsushima}によりアフィンではない。
ただし、この事実だけから判別式を除いた$\mathcal S_{N,k}$の非アフィン性や
中心元の大域選択の不存在を結論してはならない。
\end{remark}

\begin{remark}[既存結果との境界と中心元の降下]
\label{rem:pg-intrinsic-all-rank-descent-scope}
ブロック復元そのものは補題\ref{ne:blocks}に、
適合三次式の束と相対階数層の表示は既存の相対成分定理の証明に含まれる。
ここではそれらを有限\'etale被覆として整理し、
全非零階数層のgerbeの位数を同じ交換子で決めた。
ブロックの順序を選んでも、中心方向の補空間や純中心三次式の配分は選ばれない。

中心元の存在を別に確立した$H_N$-不変な零切断近傍では、
次数$d$のスカラー重みが$d$であるという既存の作用則をそのまま使える。
$\gcd(N,3)=1$なら$\theta^{-d}$によるtwistで各次数を降下できる。
全ての$N$について第三Veroneseの各次数は
\begin{equation}
 \mathcal A_{3d}\otimes\lambda^d
 \label{eq:pg-intrinsic-all-rank-third-veronese-weight}
\end{equation}
とtwistすればスカラー重み零となる。
従って共変な中心三次元とその平坦な商が構成されている場合には、
系\ref{cor:pg-intrinsic-third-veronese-descent}と同じ降下が適用できる。
この重み計算だけで、全ての$m,k$に対する同時中心元の存在や、
別の中心直線を用いた商の同型を主張するものではない。
\end{remark}

\subsection{満support層における同変中心元}
以下では$N=3k$、$k\ge2$とする。
補題\ref{lem:pg-intrinsic-block-cover}の記法
$\mathcal O_{N,k}$、$\mathcal S_{N,k}$を用いる。
この場合、$P\in\mathcal O_{N,k}$のsupportは$V$全体であり、
その三次元ブロックは順序を除いて$P$から定まる。
順序を付けた有限\'{e}tale被覆上では
\begin{equation}
 V=\bigoplus_{s=1}^k W_s,\qquad
 P=\sum_{s=1}^kP_s,\qquad
 F=\sum_{s=1}^kg_s,\qquad
 g_s\in\operatorname{Sym}^3W_s^*.
 \label{eq:pg-full-support-canonical-blocks}
\end{equation}
最後の式では$W_s^*$を他のブロックを消す$V^*$の部分空間と同一視する。
従って$g_s=F\circ\operatorname{pr}_{W_s}$であり、純中心項の配分は現れない。
$F$が滑らかであることと全$g_s$が滑らかであることは同値である。
実際、ある$g_s$の非零のgradient零点を他のブロックの零ベクトルと合わせると
$F$の非零のgradient零点となり、逆向きはgradientをブロックごとに見るだけでよい。

注意\ref{rem:pg-intrinsic-full-support-affinity}により、
$\mathcal S_{N,k}$は滑らかな既約アフィンスキームである。
具体的には、還元的安定化群
$J_{N,k}=(\mathrm{SL}_3)^k\rtimes\mathfrak S_k$を用いた
アフィン軌道$\mathrm{GL}(V)/J_{N,k}$上で、
階数$10k$の三次式のベクトル束の判別式非零部分として得られる。
以下の持上げでは、このアフィン性を用いる。

\begin{theorem}[満support層の全係数にわたる同変中心元]
\label{thm:pg-intrinsic-full-support-center}
$N=3k$、$k\ge2$とする。
$H_N=\mathrm{GL}(V)\times\mathbf G_{m,\lambda}$は
\begin{equation}
 F'(X)=\lambda^{-1}F(L^{-1}X),\qquad
 P'=\lambda(\Lambda^3L)P,\qquad
 \hbar'=\hbar
 \label{eq:pg-full-support-group-action}
\end{equation}
によって$\mathcal S_{N,k}\times\mathbb A^1_\hbar$に作用する。
普遍代数を
\begin{equation}
 \mathcal A_{F,P,\hbar}
 =\CC\langle v_1,\ldots,v_N\rangle\Big/
 \left([v_i,v_j]-\hbar\sum_aP^{ija}\sym(F_a)\right)_{i,j}
 \label{eq:pg-full-support-ambient-algebra}
\end{equation}
とする。$H_N$と$\hbar\mapsto-\hbar$で不変なZariski開
\begin{equation}
 \mathcal W_{N,k}\subset\mathcal S_{N,k}\times\mathbb A^1_\hbar,
 \qquad
 \mathcal S_{N,k}\times\{0\}\subset\mathcal W_{N,k}
 \label{eq:pg-full-support-good-open}
\end{equation}
および、その上でWeyl係数が正則な中心三次元$\Omega$が存在し、
次を満たす。
\begin{enumerate}
\item $X=Lv$による代数同型を$\Phi$とすると、
\begin{equation}
 \begin{aligned}
 \Phi(\Omega_{F',P',\hbar})&=\lambda^{-1}\Omega_{F,P,\hbar},&
 \Omega_{F,P,-\hbar}&=\Omega_{F,P,\hbar},\\
 \Omega_{F,P,\hbar}
 &=\sym\left(F-\frac{\hbar^2}{12}\mathscr H(P,F)
                  +O(\hbar^4)\right).
 \end{aligned}
 \label{eq:pg-full-support-equivariant-jet}
\end{equation}
ここで$\mathscr H(P,F)$は\eqref{eq:pg-ad-intrinsic-hessian}であり、
偶性とjetはWeyl係数についての等式である。
\item 割った交換子
\begin{equation}
 \mathcal D(p)=
 \left(\hbar^{-1}[\sym(p),v_i]\right)_{i=1}^N,
 \qquad\mathcal K=\ker\mathcal D
 \label{eq:pg-full-support-divided-kernel}
\end{equation}
を三次多項式の束から$\mathcal A_4^{\oplus N}$への写像と見る。
その核は階数$k$のベクトル束であり、核・像・余核の形成は任意の底変更と両立する。
ブロックに順序を付けた有限\'{e}tale被覆上で、$\mathcal K$は
各ブロックの中心三次直線$\mathcal K_s$の直和となる。
$\Omega$の各$\mathcal K_s$成分は零でなく、
$\mathcal J=\OO\Omega\subset\mathcal K$は$H_N$-同変な直線部分束となる。
\item $\mathcal A$と$\mathcal B=\mathcal A/(\mathcal J)$は次数別有限局所自由である。
各閉ファイバーの$\mathcal A$は次元$N$、AS parameter $N$の
Koszul AS-regular代数であり、$\Omega$は左右正則である。
各閉ファイバーの$\mathcal B$は左右injective dimension $N-1$、
AS parameter $N-3$のAS-Gorenstein代数で、
\begin{equation}
 \begin{aligned}
 \operatorname{Hilb}_{\mathcal A}(z)&=(1-z)^{-N},&
 \operatorname{Hilb}_{\mathcal B}(z)&=(1-z^3)(1-z)^{-N},\\
 \mathcal B/(\hbar)&\simeq
 \operatorname{Sym}(V^*)\otimes\OO_{\mathcal S_{N,k}}/(F).
 \end{aligned}
 \label{eq:pg-full-support-flat-quotient}
\end{equation}
特殊ファイバーの一次Poisson括弧は$\iota_{dF}P$である。
\end{enumerate}
最終的な族の底には、ブロックの順序、各ブロックの基底、
中心三次直線の相対規格化を追加しない。
中心元の一意性や全$\hbar$直線上での存在は主張しない。
\end{theorem}
\begin{proof}
\emph{古典核と局所的な量子核。}
ブロックを順序付け、その体積を$P_s$とする基底を局所的に選べば、
\eqref{eq:pg-full-support-ambient-algebra}は
\begin{equation}
 \mathcal A\simeq\bigotimes_{s=1}^k\mathcal A_{g_s,\hbar},
 \qquad
 (\mathcal K_0)_{(F,P)}=\bigoplus_{s=1}^k\CC g_s
 \label{eq:pg-full-support-tensor-and-classical-kernel}
\end{equation}
となる。テンソル積は係数底上で取る。
古典核の等号を確認する。
三次式$p$がCasimirなら各$s$について
$d_{W_s}p\wedge d g_s=0$である。
$g_s$の三つの偏微分は互いに素なので、
全変数の多項式環の分数体で得る比例係数は多項式である。
斉次次数を比較するとその係数は定数$a_s$であり、
$p-\sum_sa_sg_s$は全変数で微分が零となる。
三次斉次性からこの差は零である。
$g_s$の単項式の支持は互いに異なるので、上の和は直和となる。
この議論は補題\ref{lem:pg-det-classical-kernel}と同じUFDの議論である。
順序を替えると直和の因子が置換されるだけであり、
$\mathcal K_0$とその切断$F=\sum_sg_s$は順序を忘れた底で定義される。

各ブロックに命題\ref{prop:pg-PBW}の係数を動かした開
\eqref{eq:pg-det-universal-open}と、補題\ref{lem:pg-det-divided-kernel}を適用する。
全零切断の近傍で、各ブロック代数は次数別局所自由、
次数三・四のWeyl写像は同型、各割った交換子核$\mathcal K_s$は階数一となる。
異なるブロックの生成元は可換なので、
$\bigoplus_s\mathcal K_s$は全体の$\mathcal D$の核に入る。
テンソル積のブロック多重次数と各$\mathcal K_s$の局所分裂により、
これは$\mathcal A_3$の階数$k$の局所直和因子である。
従って$\mathcal D$は階数$r=\binom{N+2}{3}-k$の商束を経由し、
$(r+1)$次小行列式はすべて恒等的に零となる。
一方、\eqref{eq:pg-full-support-tensor-and-classical-kernel}により、
零切断の全点で階数は$r$である。
階数$r$の開に制限し、小行列式による基本変形を行うと、
$\mathcal D$は局所的に$\operatorname{diag}(I_r,0)$となる。
これで$\mathcal K$、像、余核の局所自由性と任意の底変更との両立を得る。
両辺の階数と各ファイバーでの独立性から
\begin{equation}
 \mathcal K=\bigoplus_{s=1}^k\mathcal K_s
 \label{eq:pg-full-support-quantum-kernel-sum}
\end{equation}
も成り立つ。
ここで混合単項式から追加の三次中心元が現れないことは、
この全交換子行列の階数条件によって保証されている。

これらの条件は代数と割った交換子から定まり、基底変更とブロック置換で保たれる。
従って条件の成立する開は順序を忘れた底へ降り、
全零切断を含む$H_N$-不変開$\mathcal U_{\mathrm{loc}}$を与える。
語の順序反転による反対代数との同型を用い、
$\hbar\mapsto-\hbar$に関しても不変とする。

\emph{アフィン底上の指定jetの直接持上げ。}
$\mathcal S_{N,k}$はアフィンなので、
補題\ref{lem:pg-equivariant-affine-lifting}(1)により
\begin{equation}
 D(q)\subset\mathcal U_{\mathrm{loc}},\qquad
 q\in\Gamma(\mathcal S_{N,k},\OO)[\hbar]^{H_N\times\boldsymbol\mu_2},
 \qquad q\bmod\hbar=1
 \label{eq:pg-full-support-affine-good-open}
\end{equation}
とできる。ここで$\boldsymbol\mu_2$は$\hbar$の符号と語の順序を反転する。
$D(q)$はアフィンであり、その上の$\mathcal K$は有限射影加群に対応する。
特に古典核についても$H^1(\mathcal S_{N,k},\mathcal K_0)=0$であるが、
以下では有限の厚みから直接持ち上げれば十分である。

\begin{equation}
 \Theta=F-\frac{\hbar^2}{12}\mathscr H(P,F)
 \label{eq:pg-full-support-prescribed-theta}
\end{equation}
とおく。\eqref{eq:pg-ad-intrinsic-blocks}と
命題\ref{prop:pg-universal-casimir-jet}を各ブロックに適用すると、
$\mathcal D\Theta\in\hbar^4\mathcal A_4^{\oplus N}$である。
異なるブロックは可換なので、交叉交換子は現れない。
底変更との両立から$\Theta\bmod\hbar^3$は
$\mathcal K/\hbar^3\mathcal K$の大域切断になる。
アフィン性により
\begin{equation}
 \Gamma(D(q),\mathcal K)\longrightarrow
 \Gamma\bigl(V(\hbar^3),\mathcal K/\hbar^3\mathcal K\bigr)
 \quad\text{は全射である。}
 \label{eq:pg-full-support-jet-lift-surjection}
\end{equation}
$\mathscr H(P,F)$は$F$と同じ$\lambda^{-1}$の共変性を持ち、
$\Theta$は偶である。
従って引戻しdatumを$\lambda\Phi$でtwistすると、この剰余切断は
$H_N\times\boldsymbol\mu_2$-不変となる。
補題\ref{lem:pg-equivariant-affine-lifting}(2)を閉部分$V(\hbar^3)$へ適用し、
任意の持上げをReynolds作用素で平均すれば、
$\mathcal D\Omega=0$、$\Omega\equiv\Theta\pmod{\hbar^3}$を満たす
同変かつ偶な代数的切断を得る。
偶性により差の三次係数も零なので、
\eqref{eq:pg-full-support-equivariant-jet}が成立する。
この議論は形式級数の全次数の持上げや完備化を必要としない。

順序を付けた有限\'{e}tale被覆上で、$\Omega$の各
$\mathcal K_s$成分が零になる閉集合を除く。
その和集合はブロック置換で不変なので、除いた開は元の底へ降りる。
古典極限の各成分は非零の$g_s$であるから、全零切断は残る。
この開を$\mathcal W_{N,k}$とすれば、$\Omega$は各ファイバーで非零であり、
所要の同変な直線部分束を生成する。

\emph{正則性と平坦な中心商。}
各閉ファイバーでは、$\mathcal A$は三次元Koszul AS-regular代数の
テンソル積である。
各因子のKoszul分解とAS-Gorenstein条件をテンソルする既存の議論
（定理\ref{thm:pg-ad-all-poisson}）により、
全体の次元とAS parameterはともに$3k=N$となる。

$\Omega$を\eqref{eq:pg-full-support-quantum-kernel-sum}に従って
$\Omega=\sum_s\gamma_s$と書く。
一つのブロック$s$の次数だけでfilterすると、その最高記号は
非零の$\gamma_s\in(\mathcal A_s)_3$である。
三次元因子$\mathcal A_s$はdomainなので、$\gamma_s$の左右乗法は単射である。
他の因子とのテンソル積は体上で単射性を保つ。
最高次数を比較すれば、元の$\Omega$の左右乗法も単射となる。
ここではテンソル積代数自体がdomainであるとは仮定していない。
従って各内部次数$d$について
\begin{equation}
 0\longrightarrow\mathcal J\otimes\mathcal A_{d-3}
 \longrightarrow\mathcal A_d\longrightarrow\mathcal B_d
 \longrightarrow0
 \label{eq:pg-full-support-central-quotient-sequence}
\end{equation}
は、全閉ファイバーで左の射が単射なベクトル束の列である。
最大小行列式によって局所的に分裂し、$\mathcal B_d$は有限局所自由となる。
Hilbert級数と古典極限はこの完全列から従い、
AS-Gorenstein条件は補題\ref{lem:pg-central-cubic}による。
一次Poisson括弧は\eqref{eq:pg-full-support-ambient-algebra}から直ちに従う。
\end{proof}

\begin{remark}[核と選択の範囲]
\label{rem:pg-full-support-center-scope}
この定理の$\mathcal K$は$\hbar$で割った交換子の核である。
$\hbar=0$で通常の中心は次数三全体となるため、通常の中心の形成が
任意の底変更と両立するという主張ではない。
また、各ブロックの中心直線だけから総中心直線が一意に定まるわけではない。
ここで示したのは、古典切断$F$とその二次jetを指定した一つの同変な選択の存在である。
既知の一ブロック中心元、還元群のアフィン商、アフィン底上の同変持上げを用いて、
満support層で分解の順序を保持しない同時の族を構成した。
中心変数がある場合の全supportの同変持上げは、このアフィン性の議論からは従わない。
\end{remark}

\begin{corollary}[満supportの底の次元]
\label{cor:pg-full-support-moduli-dimension}
$N=3k$、$k\ge2$の満support層について
\begin{equation}
 \dim\mathcal S_{N,k}=N^2+2k,\qquad
 \dim[\mathcal S_{N,k}/\overline H_N]=2k.
 \label{eq:pg-full-support-moduli-dimension}
\end{equation}
従って$N=6$では、このPoisson対の被約モジュライ層の次元は四である。
これは底の次元であり、固定$\hbar\ne0$での非可換代数の同型類の像の次元を主張しない。
\end{corollary}
\begin{proof}
$\dim J_{N,k}=8k$なので、$\dim\mathcal O_{N,k}=N^2-8k$である。
許される三次式の束の階数は$10k$であり、滑らかさは非空開条件である。
これより第一の等式を得る。
$\dim\overline H_N=N^2$を引けば第二の等式となる。
\end{proof}

\subsection{満support層への三次偏極での降下}
\begin{corollary}[満support層の三次Veronese族]
\label{cor:pg-full-support-veronese-descent}
$N=3k$、$k\ge2$とし、定理\ref{thm:pg-intrinsic-full-support-center}の族を取る。
全零切断を含む開部分スタック
\begin{equation}
 \mathfrak W_{N,k}=[\mathcal W_{N,k}/\overline H_N]
 \subset\mathfrak P_{n,k}^{\mathrm{red}}\times\mathbb A^1_\hbar
 \label{eq:pg-full-support-effective-open}
\end{equation}
上に、次数別有限局所自由な代数族$\mathcal A^{[3]},\mathcal B^{[3]}$が存在する。
対応するgerbeへ引き戻すと
\begin{equation}
 \mathcal A^{[3]}_d=\mathcal A_{3d}\otimes\lambda^d,
 \qquad
 \mathcal B^{[3]}_d=\mathcal B_{3d}\otimes\lambda^d
 \quad(d\ge0)
 \label{eq:pg-full-support-third-veronese}
\end{equation}
であり、次数一の大域的な左右正則中心元$\omega$について
\begin{equation}
 \mathcal B^{[3]}=\mathcal A^{[3]}/(\omega)
 \label{eq:pg-full-support-veronese-quotient}
\end{equation}
となる。$\hbar=0$での相対$\operatorname{Proj}$は普遍の周囲射影空間族と
三次超曲面族であり、新しい次数一の偏極は各幾何学的ファイバーで$\OO(3)$である。
他方、全層のスカラーgerbeは位数三を持ち、元の次数一の線形化を
一つの可逆層の次数冪で補正して降下させることはできない。
\end{corollary}
\begin{proof}
$K_{\mathrm{sc}}$は$\mathcal A_d,\mathcal B_d$に重み$d$、
指標$\lambda$に重み$-3$で作用する。
従って\eqref{eq:pg-full-support-third-veronese}の各成分の重みは零であり、
乗法と単位もgerbeの重み零の降下
\cite[Proposition~4.10]{PBerghSchnurerGerbes}によって降りる。
中心元の共変性\eqref{eq:pg-full-support-equivariant-jet}は
$\mathcal J\otimes\lambda\simeq\OO$という同変な自明化を与えるので、
$\Omega$は新しい次数一の大域切断$\omega$を定める。
式\eqref{eq:pg-full-support-central-quotient-sequence}の$d$を$3d$に替え、
$\lambda^d$で捻ると、左端は
\begin{equation}
 (\mathcal J\otimes\mathcal A_{3d-3})\otimes\lambda^d
   \simeq\mathcal A_{3d-3}\otimes\lambda^{d-1}
   =\mathcal A^{[3]}_{d-1}
 \label{eq:pg-full-support-veronese-exact-sequence}
\end{equation}
となる。よって商の同定、正則性、次数別局所自由性が従う。
古典極限では通常の三次Veroneseと底の可逆層の次数冪による捻りを取っている。
いずれも相対$\operatorname{Proj}$を保ち、偏極は各ファイバー上で$\OO(3)$となる。
この同定は自明化した底上の通常の同定から降下するので、
周囲族が大域的にベクトル束の射影化であるとは仮定しない。
最後の非存在は命題\ref{prop:pg-intrinsic-all-rank-gerbe}から従う。
\end{proof}

\begin{remark}[三次四様体と比較の範囲]
\label{rem:pg-full-support-fourfold-scope}
$N=6$、$k=2$では、正規形定理により三次四様体の一般階数四のPoisson対を
全て扱っている。底の次元は系\ref{cor:pg-full-support-moduli-dimension}の通り四であり、
その全体で二つの三生成元量子代数と総中心三次商を同時に構成し、
三次偏極によってPoisson対のモジュライへ降ろしたことになる。
この主張は、固定非零$\hbar$での同型類像の次元や一意性を含まない。

三次Veronese自体のAS-regular性は主張しない。
左右noetherianな幾何学的ファイバーに限れば、
注意\ref{rem:pg-intrinsic-veronese-qgr}の直接証明と
Smithの既知のVeronese同値\cite{PSmithVeronese}をそのまま適用できる。
これは相対的な$qgr$の降下や、全ファイバーの滑らかさの証明ではない。
中心変数を持つ複数ブロックの全係数空間への同変中心元の延長と、
異なる中心直線による商の抽象的同型・Morita・$qgr$比較は別に残る。
\end{remark}
